\documentclass[11pt,a4paper]{article}

%% PACHETE MATEMATICE
\usepackage{amsmath,amssymb,amsfonts}
\usepackage{amsthm}
\usepackage{mathpazo}
\usepackage{eulervm}
\usepackage{enumerate}
\usepackage{color}
\usepackage{graphicx}
\usepackage[all]{xy}
\usepackage{xcolor}
\usepackage{framed}
\usepackage[unicode]{hyperref}
\setcounter{MaxMatrixCols}{20}
\usepackage{multicol}
%\usepackage[romanian]{babel}


%% ASPECT
\linespread{1.05}
\usepackage[T1]{fontenc}
\usepackage[utf8x]{inputenc}
\usepackage[margin=2cm]{geometry}
% \usepackage[color]{showkeys}
\usepackage{anyfontsize}
\usepackage{titlesec}
\titleformat*{\section}{\large\bfseries}

\usepackage{fancyhdr}
\pagestyle{fancy}
\rhead{\small \color{gray} Notițe de Adrian Manea}
\lhead{\small \color{gray} BCD-Aero, 2017---2018, L. Costache \& M. Olteanu}


\usepackage[autostyle = false, style = german]{csquotes}
\usepackage[romanian]{babel}
\newcommand\qq{\enquote}

%% COMENZILE MELE
\renewcommand*{\proofname}{Dem.:}
\newcommand\bb{\mathbb}
\newcommand\dr{\mathrm}
\newcommand\kal{\mathcal}
\newcommand\fk{\mathfrak}
\newcommand\op{\oplus}
\newcommand\ot{\otimes}
\newcommand\ds{\displaystyle}
\newcommand\id{\indent}
\newcommand\nid{\noindent}
\newcommand\seq{\subseteq}
\newcommand\sneq{\subsetneq}
\newcommand\speq{\supseteq}
\newcommand\spneq{\supsetneq}
\newcommand\rar{\rightarrow}
\newcommand\Rar{\Rightarrow}
\newcommand\xrar{\xrightarrow}
\newcommand\lar{\leftarrow}
\newcommand\Lar{\Leftarrow}
\newcommand\xlar{\xleftarrow}
\newcommand\lrar{\leftrightarrow}
\newcommand\Lrar{\Leftrightarrow}
\newcommand\ideal{\trianglelefteq}
\newcommand\ai{\text{ a.\^{i}. }}
\newcommand\sk{(\textit{Schi\c{t}\u{a}}) }
\newcommand\rhup{\rightharpoonup}
\newcommand\rhdn{\rightharpoondown}
\newcommand\lhup{\leftharpoonup}
\newcommand\lhdn{\leftharpoondown}
\newcommand\vid{\emptyset}
\newcommand\wh{\widehat}
\newcommand\ol{\overline}
\newcommand\NN{\mathbb{N}}
\newcommand\RR{\mathbb{R}}
\newcommand\ZZ{\mathbb{Z}}
\newcommand\QQ{\mathbb{Q}}
\newcommand\CC{\mathbb{C}}

%% STILURILE MELE
\newtheoremstyle{dotless}{}{}{\itshape}{}{\bfseries}{:}{ }{}

\theoremstyle{dotless}
\newtheorem{theorem}{Teorem\u{a}}[section]
\newtheorem{proposition}{Propozi\c{t}ie}[section]
\newtheorem{lemma}{Lem\u{a}}[section]
\newtheorem{corollary}{Corolar}[section]
\newtheorem{exercise}{Exerci\c{t}iu}

\newtheoremstyle{dotlessdr}{}{}{}{}{\bfseries}{:}{ }{}
\theoremstyle{dotlessdr}
\newtheorem{example}{Exemplu}[section]
\newtheorem{definition}{Defini\c{t}ie}[section]
\newtheorem{remark}{Observa\c{t}ie}[section]




\begin{document}

\begin{center}
  {\large\textbf{Seminar 6 \\ Mulțimi numărabile și nenumărabile}}
\end{center}

\vspace{1cm}

\section{Aspecte teoretice}
\label{sec:teorie-multimi}



Din punct de vedere al cardinalului (numărului de elemente), mulțimile pot fi
\emph{finite} sau \emph{infinite}. Mulțimile infinite, la rîndul lor, se împart în
mulțimi \emph{numărabile} și mulțimi \emph{nenumărabile}.

\begin{definition}\label{def:numarabil}
  Intuitiv, o mulțime se numește \emph{numărabilă} dacă elementele sale pot fi
  enumerate.

  Formal, o mulțime $X$ se numește \emph{numărabilă} dacă există o bijecție
  $f : X \to \NN$, unde $\NN$ este mulțimea numerelor naturale.

  O mulțime se numește \emph{cel mult numărabilă} dacă este finită sau numărabilă.
\end{definition}

Pentru o mulțime infinită (numărabilă sau nu), se poate demonstra următorul rezultat:

\begin{theorem}\label{thm:num-subm}
  O mulțime este infinită dacă și numai dacă există o bijecție între ea și o submulțime
  proprie a ei.
\end{theorem}

De exemplu, între $\NN$ și mulțimea numerelor naturale pare există bijecția $f(x) = 2x$,
ceea ce ne arată că $\NN$ este infinită.

Vom nota cardinalul mulțimii numerelor naturale cu $\aleph_0$ (\emph{alef} zero).

Se poate arăta că: $| \NN | = | \ZZ | = | \QQ | = \aleph_0$, dar $| \RR | \neq \aleph_0$.

Mai mult, folosind \emph{argumentul diagonal al lui Cantor} (v. cursul sau
alte resurse), putem arăta că $\RR$ este \emph{nenumărabilă}. Și mai mult, chiar
se poate arăta (ceva mai greu) că $| \RR | = 2^{\aleph_0}$.

Următorul rezultat ne arată cum putem decide dacă o mulțime este numărabilă sau nu:

\begin{theorem}\label{thm:num}
  Fie $X$ și $Y$ două mulțimi nevide și fie funcția $f : X \to Y$.
  \begin{enumerate}[(a)]
  \item Dacă $X$ și $Y$ sînt numărabile, atunci $X \cup Y$ este numărabilă;
  \item Mai general, o reuniune numărabilă de mulțimi numărabile este numărabilă;
  \item Dacă $f$ este injectivă, iar $Y$ este numărabilă, atunci $X$ este numărabilă;
  \item Mai general, dacă $f$ este injectivă, atunci $|X| \leq |Y|$;
  \item Dacă $f$ este surjectivă, iar $X$ este numărabilă, atunci $Y$ este
    numărabilă;
  \item Mai general, dacă $f$ este surjectivă, atunci $|X| \geq |Y|$;
  \item Dacă $f$ este bijectivă, atunci $| X | = | Y |$.
  \item Produsul cartezian $\NN \times \NN$ este numărabil;
  \item Mai general, un produs cartezian \emph{finit} de mulțimi numărabile este
    numărabil;
  \item Un produs cartezian infinit de mulțimi (cel mult) numărabile
    este \emph{nenumărabil} (v. argumentul lui Cantor, varianta cu $\{ 0, 1 \}$).
  \end{enumerate}
\end{theorem}

Așadar, pentru a verifica dacă o mulțime este numărabilă, trebuie să judecăm
cu ajutorul funcțiilor.


%%%%%%%%%%%%%%%%%%%%%%%%%%%%%%%%%%%%%%%%%%%%%%%%%%%%%%%%%%%%%%%%%%%%%%

\section{Exerciții}
\label{sec:exc}

Decideți care din următoarele mulțimi sînt numărabile:
\begin{enumerate}[(a)]
\item $A = \{ \frac{m^2 + n^2}{2} \mid m, n \in \NN \}$;
\item $B = \{ \frac{x^2 + i\sqrt{x}}{2} \mid x \in [0, \infty) \cap \QQ \}$;
\item $C = \{ f : \NN \to \{ 0, 1 \} \}$;
\item $D = \{ f : \{0, 1, \dots, 9 \} \to \QQ \}$;
\item $E = \{ (b_0, b_1, \dots) \mid b_i \in \{ x, y, z \}, \forall i \}$;
\item $F = \{ it \mid t \in \RR - \QQ \}$.
\item $G = \{ \frac{n + \sqrt{n^2 + 1}}{2} \mid n \in \QQ \cap [0, \infty) \}$;
\item $H = \{ \frac{1 - i\sqrt{n}}{2} \mid n \in \NN^\ast \}$;
\item $I = \{ \frac{m^2 + \sqrt{n}}{2} \mid m, n \in \QQ_+ \}$;
\item $J = \{ n \in \ZZ \mid 5 \nmid n, 7 \mid n \}$;
\item $K = \{ \frac{m^2 + i\sqrt{3}}{n^2 + 1} \mid m, n \in \ZZ \}$;
\end{enumerate}

\end{document}
